\documentclass[a4paper]{article}

\usepackage{anyfontsize}
\usepackage{amsmath}

\usepackage[no-math]{fontspec}
\setmainfont{XCharter}
\usepackage{unicode-math}
\unimathsetup{bold-style=ISO,sans-style=italic}
\setmathfont{XCharter-Math.otf}
\AtBeginDocument{
  \let\uglymathbf\mathbf
  \renewcommand\mathbf\symbf
  \let\uglymathsf\mathsf
  \renewcommand\mathsf\symsf
}

\AtBeginDocument{
  \let\uglyepsilon\epsilon
  \renewcommand\epsilon\varepsilon
}

\usepackage{indentfirst}
\setlength{\parindent}{0em}
\usepackage{autobreak}
\allowdisplaybreaks
\usepackage{parskip}
\usepackage{geometry}
\geometry{top=3cm,bottom=3cm,left=2cm,right=2cm}

\usepackage[fontset=none]{ctex}
\setCJKmainfont{Adobe Song Std}[BoldFont=Noto Sans CJK SC Medium]

\begin{document}

\title{\textbf{天体光谱学作业}}
\author{1900011604\ \ 张植竣\ \ 北京大学}
\date{2022年4月1日}
\maketitle

\begin{enumerate}

    \item[1.]
    由于红移$z=\dfrac{\lambda - \lambda_0}{\lambda_0}$，其中$\lambda$为观测到的波长，$\lambda_0$为静止波长，可以计算四条波长的红移分别为：
    \[
        z_1 = 3.34008 \times 10^{-3},\quad
        z_2 = 3.34063 \times 10^{-3},\quad
        z_3 = 3.33244 \times 10^{-3},\quad
        z_4 = 3.33699 \times 10^{-3}
    \]
    可见它们的红移大致相同，说明这些发射线有相同的Doppler移动，红移值平均为：
    \[
        z = \frac{1}{4}(z_1 + z_2 + z_3 + z_4) = 3.3375 \times 10^{-3}
    \]
    这个星系与我们的视向速度为：
    \[
        v = cz = 1 \times 10^6\,\mathrm{m/s}
    \]
    由于$z>0$，星系在远离我们。

    \bigskip
    \item[2.]
    经过测量估计，$\mathrm{H}40\alpha$的视向速度为：$v_\alpha = 53.00\,\mathrm{km/s}$，$\mathrm{H}50\beta$的视向速度为：$v_\beta = 656.56\,\mathrm{km/s}$，则两条谱线的红移分别为：
    \[
        z_\alpha = \frac{v_\alpha}{c} = 1.77 \times 10^{-4},\quad
        z_\beta = \frac{v_\beta}{c} = 2.19 \times 10^{-3}
    \]
    对于频率，红移$z = \dfrac{f_0-f}{f}$，其中$f$为观测到的频率，$f_0$为静止频率，则观测到$\mathrm{H}50\beta$的波长为：
    \[
        f_\beta = \frac{f_{\alpha 0}}{1 + z_\beta} = 98.8066\,\mathrm{GHz}
    \]
    而$f_{\beta0}$的真实红移应该与$f_{\alpha0}$相同，即$z_\alpha=1.77\times10^{-7}$，于是有：
    \[
        f_{\beta0} = f_\beta(1 + z_\alpha) = f_{\alpha0}\frac{c + v_\alpha}{c + v_\beta} = 98.824\,\mathrm{GHz}
    \]

    \bigskip
    \item[3.]
    氢原子的能级：
    \[
        E_n = \frac{E_1}{n^2},\quad
        n = 1, 2, 3, \cdots
    \]
    而对于氢原子，Rydberg公式：
    \[
        \frac{1}{\lambda} = R_\mathrm{H}\left(\frac{1}{n_f^2} - \frac{1}{n_i^2}\right),\quad
        R_\mathrm{H} = \frac{\mu}{4\pi c\hbar^3}\left(\frac{e^2}{4\pi\epsilon_0}\right)^{\!2}
    \]
    再与$E=\dfrac{hc}{\lambda}$联立，得用氢原子Rydberg常数表达氢原子的能级为：
    \[
        E_n = \frac{hcR_\mathrm{H}}{n^2},\quad
        n = 1, 2, 3, \cdots
    \]
    $R_\infty$与$R_\mathrm{H}$的区别在于：
    \[
        \left\{
            \begin{aligned}
                & R_\infty = \frac{m_e}{4\pi c\hbar^3}\left(\frac{e^2}{4\pi\epsilon_0}\right)^{\!2} \\
                & R_\mathrm{H} = \frac{\mu}{4\pi c\hbar^3}\left(\frac{e^2}{4\pi\epsilon_0}\right)^{\!2},\quad
                \mu = \frac{M m_e}{M + m_e} = \dfrac{m_e}{1 + \dfrac{m_e}{M}}
            \end{aligned}
        \right.
    \]
    此处$m_e$为电子质量，$M$为原子核质量，$\mu$为约化质量。

    对于重元素类氢离子，$M_Z \gg m_e$，$\dfrac{m_e}{M_Z} \ll 1$，则$m_e$相比$\mu_\mathrm{H}$与$\mu_Z$更为接近，证明如下：
    \[
        |m_e - \mu_Z| = \left|\frac{m_e(M_Z + m_e)}{M_Z + m_e} - \frac{M_Z m_e}{M_Z + m_e}\right| = \frac{m_e^2}{M_Z + m_e} = \frac{m_e^2(M_H + m_e)}{(M_Z + m_e)(M_H + m_e)}
    \]
    \[
        |\mu_H - \mu_Z| = \left|\frac{M_H m_e}{M_H + m_e} - \frac{M_Z m_e}{M_Z + m_e}\right| = \frac{m_e^2(M_Z - M_H)}{(M_Z + m_e)(M_H + m_e)}
    \]
    又因为：
    \[
    M_H + m_e \approx M_H \ll M_Z - M_H \approx M_Z
    \]
    可见用Rydberg常数$R_\infty$比用氢原子Rydberg常数$R_\mathrm{H}$更合适。

    对于铁的类氢离子$\mathrm{Fe}^{25+}$，$Z = 26$，其$\mathrm{Ly}\alpha$的波数为：
    \[
        \frac{1}{\lambda} = R_\mathrm{Fe}\left(\frac{1}{1^2}-\frac{1}{2^2}\right) \approx Z^2 R_\infty\left(\frac{1}{1^2}-\frac{1}{2^2}\right) = 5.56368 \times 10^9\,\mathrm{m^{-1}}
    \]

    \bigskip
    \item[4.]
    由电偶极跃迁的选择定则，$4\mathrm{p}$态可以跃迁的路径如下：
    4p-3s, 4p-3d, 4p-2s, 4p-1s, 3s-2p, 3d-2p, 2p-1s

    则它可以有6种波长：$R=1.0967758\times10^{7}\,\mathrm{m}^{-1}$
    \[
        \lambda_{4 \to 3} = 1.87563\,\mu\mathrm{m},\quad
        \lambda_{4 \to 2} = 0.48627\,\mu\mathrm{m},\quad
        \lambda_{4 \to 1} = 97.25476\,\mathrm{nm}
    \]
    \[
        \lambda_{3 \to 2} = 0.65647\,\mu\mathrm{m},\quad
        \lambda_{3 \to 1} = 0.10257\,\mu\mathrm{m},\quad
        \lambda_{2 \to 1} = 121.56845\,\mathrm{nm}
    \]

    由电偶极跃迁的选择定则，$4\mathrm{s}$态可以跃迁的路径如下：
    4s-3p, 4s-2p, 3p-2s, 3p-1s, 2p-1s

    则它可以有5种波长：
    \[
        \lambda_{4 \to 3} = 1.87563\,\mu\mathrm{m},\quad
        \lambda_{4 \to 2} = 0.48627\,\mu\mathrm{m}
    \]
    \[
        \lambda_{3 \to 2} = 0.65647\,\mu\mathrm{m},\quad
        \lambda_{3 \to 1} = 0.10257\,\mu\mathrm{m},\quad
        \lambda_{2 \to 1} = 121.56845\,\mathrm{nm}
    \]

    \bigskip
    \item[5.]
    对于氘原子，由于$M_\mathrm{D} = 3670.4m_e$，约化质量$\mu = \dfrac{M_\mathrm{D} m_e}{M_\mathrm{D} + m_e}=1.00027245m_e$，则Rydberg常数为：
    \[
        R_\mathrm{D} = \frac{\mu}{4\pi c\hbar^3}\left(\frac{e^2}{4\pi\epsilon_0}\right)^{\!2} = 1.0970743\times10^{7}\,\mathrm{m}^{-1}
    \]
    则氘原子的$\mathrm{H}\alpha$波长为：
    \[
        \lambda_\mathrm{D} = R_\mathrm{D}^{-1}\left(\frac{1}{2^2} - \frac{1}{3^2}\right)^{-1} = 6562.91034\,\text{\AA}
    \]
    同理，氢原子（$M_\mathrm{H}=1836.1m_e$）的$\mathrm{H}\alpha$波长为：
    \[
        \lambda_\mathrm{H} = R_\mathrm{H}^{-1}\left(\frac{1}{2^2} - \frac{1}{3^2}\right)^{-1} = 6564.69617\,\text{\AA}
    \]
    分辨这两根谱线需要的光谱分辨率为：
    \[
        R = \frac{\lambda}{\Delta\lambda} = \frac{\lambda_\mathrm{H}}{|\lambda_\mathrm{H} - \lambda_\mathrm{D}|} = 3675.99165 = 81.55417\,\mathrm{km/s}
    \]
    对于$\mathrm{H}\alpha$和$\mathrm{Ly}\alpha$，所需要的分辨率是一样的，证明如下：
    \[
        R = \frac{\lambda}{\Delta\lambda} = \frac{\lambda_\mathrm{H}}{|\lambda_\mathrm{H} - \lambda_\mathrm{D}|} = \frac{R_\mathrm{H}^{-1}}{|R_\mathrm{H}^{-1} - R_\mathrm{D}^{-1}|} = \frac{R_\mathrm{D}}{|R_\mathrm{D} - R_\mathrm{H}|} = \frac{\mu_\mathrm{D}}{|\mu_\mathrm{D} - \mu_\mathrm{H}|}
    \]
    这个比值只和两种原子的原子核质量有关，与电子的能级无关。

    \bigskip
    \item[6.]
    氢原子的最高能级取决于其原子半径$R_n$的大小，联立：（$\mu=\dfrac{M m_e}{M + m_e}$为约化质量）
    \[
        \mu\omega r^2 = n\hbar,\quad
        \mu\omega^2 r = \frac{e^2}{4\pi\epsilon_0 r}
    \]
    得氢原子半径：
    \[
        R_n = n^2 a,\quad
        n \in \mathbb{N}^+
    \]
    其中$a = \dfrac{4\pi\epsilon_0 \hbar^2}{\mu e^2}$为氢原子的Bohr半径，约为$0.5294654\,\text{\AA}$。

    则在恒星中：数密度$N = 10^{16}\,\mathrm{cm^{-3}}$，则最大半径为：
    \[
        r = \sqrt[3]{\frac{3}{4\pi N}} = 287.9412\,\text{\AA}
    \]
    解不等式$r \geqslant n^2 a$得$n$最大可以取的正整数值为$n=23$，即氢原子最高能级为$23$。

    同理，对于$\mathrm{H\,II}$，最大半径为：
    \[
        r' = \sqrt[3]{\frac{3}{4\pi N}} = 287.9412\,\mu\mathrm{m}
    \]
    解不等式$r' \geqslant n^2 a$得$n$最大可以取的正整数值为$n=2332$，即氢原子最高能级为$2332$。

\end{enumerate}

\end{document}
